% page 583 du fac-simile (image p-582 du scan)
% bloc 167.6 x 242.0 mm ; marge gauche 34.4 mm
\pgc{12.700mm}{0.000mm}{
\l{\cel{53}575}
\l{}
\l{\sou{tenera}.\cel{2}Qua havas por ulu afeciono qua abundas ye sentemeso e}
\l{\cel{3}dolceso. - EFIS.}
\l{}
\l{\sou{tenio}. Vermo rimeno-forma, qua kreskas en la intestini di la}
\l{\cel{3}homo e di ula animali. - L.\cel{1}\sou{taenia}.\cel{2}EFIS.}
\l{}
\l{\sou{teniso}. Sporto quan onu praktikas per raketi e baloneti, sur}
\l{\cel{3}tereno equipita specala. - DEFIS.}
\l{}
\l{\sou{tenono}. (\sou{tekn.}) Parto salianta, situita ye la extremajo di peco}
\l{\cel{3}de ligno, de metalo, de marmoro, e c., por ajusteso a parto,}
\l{\cel{3}a mortezo, korespondanta en altra peco qua esas asemblenda an}
\l{\cel{3}la unesma.- EF.}
\l{}
\l{\sou{tenoro}. I. (\sou{muziko).}\cel{1}Voco virala qua iras de la unesma \textquotedbl{}do\textquotedbl{} di}
\l{\cel{3}la alto, til la \textquotedbl{}sol\textquotedbl{} duesma di la violino. - II. Kantisto qua}
\l{\cel{3}havas tala voco. - DEFIRS.}
\l{}
\l{\sou{tensar}. (\sou{trans}.)I. Igar rigida (korpo plu o min elastika) per}
\l{\cel{3}igar lu plu longa, per tiro de un ek lua extremaji a la altra.}
\l{\cel{3}- II. (\sou{pri horlojo murala)}. Ri-suprigar la pezo-bloki kande}
\l{\cel{3}li ter-venis. - EFIS.}
\l{}
\l{\sou{tentar.}\cel{1}(\sou{trans.}) Solicitar (ulu) ad agar ulo interdiktita,}
\l{\cel{3}per olua atrakto. - EFIS.}
\l{}
\l{\sou{tentakulo}\cel{1}. (\sou{zool.}) Apendico movebla ye qua esas provizita ula}
\l{\cel{3}animali, e qua uzesas da li kom, precipue, organo di tusho-}
\l{\cel{3}senso. EFIS.}
\l{}
\l{\sou{tenua}. Preske neprenebla pro lua mikregeso, dinegeso, delikat-}
\l{\cel{3}eso. - EFIS.}
\l{}
\l{\sou{teodiceo}. (\sou{filozof}.)I.Traktato pri la yusteso di deo. - II. Parto}
\l{\cel{3}di filozofio, qua traktas la existo e la atributi di deo. -}
\l{\cel{3}DEFIRS.}
\l{}
\l{\sou{teodolito}. (\sou{astron. e geom.})\cel{2}Instrumento di geodezio, quan onu}
\l{\cel{3}uzas por relevar mapi, por mezurar la anguli reduktita a la}
\l{\cel{3}horizonto, la disti zenitala e laazimuti. - DEFIR.}
\l{}
\l{\sou{teogo}nio. Genito di la dei ; ensemblo de deaji di qui la kulto}
\l{\cel{3}formacas la sistemo religiala di populo politeista. - DEFIRS.}
\l{}
\l{\sou{teokrato.}\cel{2}Singla ek la chefi di teokratio.- DEFS.}
\l{}
\l{\sou{teokratio.}\cel{2}Formo di guvernado segun qua la chefi di la naciono,}
\l{\cel{3}de la kasto \textquotedbl{}sacerdoti\textquotedbl{}, egardesas kom ministri di deo.- DEFIRS}
\l{}
\l{\sou{teologala.}\cel{1}Dicesas pri singla de la tri vertui di qui deo esas}
\l{\cel{3}la objekto ed un ek lua atributi kom motivo (kredo, espero e}
\l{\cel{3}karitato). - FIS.}
\l{}
\l{teologio. Doktrino religiala pri la kozi deala. - DEFIRS.}
}
